\section{The multiplicative core: adequacy and complexity}
\label{sec:p08-imll-complete}

The result concerns exactly \hyperref[sheet:imll]{the registered IMLL calculus}: one succedent, finite
multiset antecedents, and the connectives \(\mathbf1,\otimes,\multimap\).
It does not cover the subsequent exponential extension.
The algebraic interpretation and empty-context convention are those of
\cref{sec:relevance-linear-03,def:p08-imll-empty-tensor}.

\begin{hbtheorem}{thm:p08-imll-historical-sketch}{sketched=annex:p8.imll.source}{Source claim}
The inherited result at
\hyperref[thm:p08-source-imll-adequacy]{Soundness, completeness, and cut elimination}
announces adequacy for commutative residuated partially ordered monoids
and admissibility of cut. Its adjacent argument is a sketch.
The following theorem supplies the complete argument for the same
multiplicative calculus, separately from the historical sketch.
\end{hbtheorem}

\begin{hbtheorem}{thm:p08-imll-complete}{proved=annex:p08-imll-complete}{Soundness completeness and cut elimination for IMLL}
For a finite multiset \(\Gamma\) and a formula \(A\) of
\hyperref[sheet:imll]{the registered IMLL calculus}, the following are equivalent:
\begin{enumerate}
 \item \(\Gamma\Rightarrow A\) has a derivation in that calculus.
 \item \(\Gamma\Rightarrow A\) has a cut-free derivation.
 \item In every commutative residuated partially ordered monoid and
       every valuation, the product of the values in \(\Gamma\)
       is below the value of \(A\).
\end{enumerate}
The product of an empty antecedent is the monoid unit.
No weakening, contraction, distributive lattice structure, or
exponential operation is assumed.
\end{hbtheorem}

\begin{hbmethod}{def:p08-imll-cut-procedure}{Eliminating a cut}
First eliminate cuts in the immediate subderivations. For a remaining
cut with cut-free premises, mark the particular occurrence of its
cut formula in the right premise. If an identity ends either premise,
remove the cut directly. Otherwise commute the cut above a last rule
that does not introduce its marked formula. If both last rules
introduce that formula, replace a tensor cut by two cuts on its factors,
an implication cut by cuts on its domain and codomain, and a unit cut
by the right premise before unit insertion. The proof establishes all
these cases. A lexicographic induction on formula size and the sum of
premise heights proves termination. The procedure does not copy or
discard resources; equal formula occurrences remain separately marked.
\end{hbmethod}

The canonical algebra in the proof is a countermodel method as well:
if a sequent is not derivable, the quotient of formulas by mutual
derivability, with atoms interpreted by their own classes, falsifies it.
This is an existence construction, not a claim of an efficient finite
countermodel algorithm. The inherited algebraic and proof-theoretic
setting is discussed in \cite{Girard1987,GalatosJipsenKowalskiOno2007}.

\begin{hbtheorem}{thm:p08-imll-proof-bound}{proved=annex:p08-imll-proof-bound}{Finite proof search and polynomial certificates}
Let \(N\) be the total number of atom, connective and unit occurrences
in an input sequent of \hyperref[sheet:imll]{the registered IMLL calculus}. Every cut-free proof of it
has at most \(N\) nodes. Exhaustive backward search over the registered
cut-free rules terminates and decides derivability. With an effective
finite-string encoding of atoms and formulas, a successful proof has
polynomial encoded size and is checkable in polynomial time.
Consequently derivability belongs to \(\mathsf{NP}\).
This proves membership, not the separate NP-hardness assertion.
\end{hbtheorem}

\subsection{The lower complexity bound}

\begin{hbdefinition}{def:p08-imll-partition-encoding}{The partition encoding}
For a formula \(A\) and a positive integer \(n\), let \(A^{\otimes n}\)
denote a fixed parenthesization of \(n\) copies of \(A\), not an additional
connective.  A unary 3-PARTITION instance consists of positive integers
\(m,B,a_1,\ldots,a_{3m}\), with
\[
 \sum_{i=1}^{3m}a_i=mB,\qquad B/4<a_i<B/2,
\]
written in unary.  Its question is whether the indices can be partitioned
into \(m\) triples, each having sum \(B\).  For distinct atoms \(p,q\), put
\[
 F_i=p^{\otimes a_i}\multimap q,\qquad
 D=p^{\otimes B}\multimap q^{\otimes3}.
\]
The encoded sequent is
\[
 F_1,\ldots,F_{3m}\Rightarrow D^{\otimes m}.
\]
Repeated numerical values still give separate antecedent occurrences.
\end{hbdefinition}

\begin{hbtheorem}{thm:p08-imll-partition}{proved=annex:p08-imll-partition}{Correctness of the encoding}
The sequent in \cref{def:p08-imll-partition-encoding} is derivable in
\hyperref[sheet:imll]{the registered IMLL calculus} if and only if its
3-PARTITION instance has a solution.  The translation has length
\(O(mB+m)\) and is polynomial-time for unary input.
\end{hbtheorem}

\begin{hbcorollary}{thm:p08-imll-np-complete}{proved=annex:p08-imll-np-complete}{NP-completeness}
Derivability in the registered IMLL calculus is NP-complete under
polynomial-time many-one reductions.  Hardness already holds for the
two-atom, unit-free sequents of \cref{def:p08-imll-partition-encoding}.
\end{hbcorollary}

The hardness premise is the strong NP-completeness of 3-PARTITION,
equivalently the NP-completeness of its unary version; its source is
\textcite[Theorem~4.4]{GareyJohnson1979}.  The encoding and its converse
are proved here.  Unary encoding matters: expanding a binary integer
into that many tensor factors would not in general be polynomial.
